\documentclass[11pt]{article}
\usepackage[margin=1.1in]{geometry}
\usepackage{amsmath,amssymb,amsthm}
% Optional packages: the file also compiles on a minimal TeX installation without them.
\IfFileExists{booktabs.sty}{\usepackage{booktabs}}{\newcommand{\toprule}{\hline}\newcommand{\midrule}{\hline}\newcommand{\bottomrule}{\hline}}
\IfFileExists{array.sty}{\usepackage{array}}{}
\IfFileExists{longtable.sty}{\usepackage{longtable}}{\newenvironment{longtable}[1]{\begin{center}\begin{tabular}{#1}}{\end{tabular}\end{center}}\let\endhead\relax}
\IfFileExists{hyperref.sty}{\usepackage[hidelinks]{hyperref}}{}

\newtheorem{theorem}{Theorem}
\newtheorem{proposition}[theorem]{Proposition}
\newtheorem{lemma}[theorem]{Lemma}
\newtheorem{corollary}[theorem]{Corollary}
\theoremstyle{remark}
\newtheorem{remark}[theorem]{Remark}
\DeclareMathOperator{\Aut}{Aut}
\newcommand{\C}{\mathbb{C}}

\title{Fusions of the point--distance schemes of generalized hexagons and octagons,\\
and the orbital algebra of the Tits group on $1755$ points}
\author{Draft --- authors to be added}
\date{28 October 2026}

\begin{document}
\maketitle

\begin{abstract}
Let $T={}^2F_4(2)'$ be the Tits group and $H=2.2^{8}{:}5{:}4$. The action of $T$ on $\Omega=T/H$, $|\Omega|=1755$, has rank~$5$; its orbitals are the distance relations of the collinearity graph of the Ree--Tits generalized octagon of order $(2,4)$, a distance-regular graph with intersection array $\{10,8,8,8;1,1,1,5\}$. We record the complete orbital algebra (all $125$ intersection numbers, the eigenmatrix, multiplicities and Krein parameters) and determine all fusions of its orbitals. More generally, for the point--distance scheme of a generalized $2d$-gon of order $(s,t)$ with $d\in\{3,4\}$ and integers $s,t\ge 2$, we show that a proper fusion scheme exists if and only if $s=t$, and that it is then unique. For hexagons ($d=3$) the fusion is strongly regular with the parameters of the collinearity graph of $Q(6,q)$; for octagons ($d=4$) the case $s=t$ is excluded by the Feit--Higman condition that $2st$ be a square. Consequently no strongly regular graph, and no association scheme other than the trivial ones, arises by merging orbitals of $T$ on $\Omega$. The proofs are elementary, rest on a sign analysis of differences of eigenvalue sums, and are backed by exact computer algebra and an exhaustive test of $841$ parameter pairs for each of $d=3,4$.
\end{abstract}

\section{Introduction}

The Tits group $T={}^2F_4(2)'$ has order $17\,971\,200=2^{11}\cdot3^3\cdot5^2\cdot13$. It has a subgroup $H\cong 2.2^8{:}5{:}4$ of order $10\,240$, hence acts transitively on $\Omega=T/H$ with $|\Omega|=1755$. The permutation representation of degree $1755$ is listed in the ATLAS of Group Representations~\cite{ATLASrep} and has rank $5$ with subdegrees $1,10,80,640,1024$. Let $R_0,\dots,R_4$ be the orbitals (with $R_0$ the diagonal) and $A_i$ their adjacency matrices; then $A_iA_j=\sum_{k}p_{ij}^kA_k$ with non-negative integers $p_{ij}^k$ (intersection numbers).

\paragraph{Results.}
\begin{theorem}[Orbital algebra]\label{thm:A}
All five orbitals are symmetric, so the orbital algebra is a commutative $5$-class association scheme, and $R_i$ is the distance-$i$ relation of the distance-regular graph $\Gamma=(\Omega,R_1)$ with intersection array $\{10,8,8,8;1,1,1,5\}$. The intersection matrices $B_i=(p_{ij}^k)_{j,k}$ are given in \eqref{eq:B}, the eigenmatrix and multiplicities in \eqref{eq:P}.
\end{theorem}

\begin{theorem}[Fusions of generalized hexagons and octagons]\label{thm:B}
Let $d\in\{3,4\}$ and let $\Gamma$ be a distance-regular graph with the intersection array of a generalized $2d$-gon of order $(s,t)$, where $s,t\ge2$ are integers. The distance scheme $\{A_0,\dots,A_d\}$ has a proper fusion scheme if and only if $s=t$. In that case the fusion is unique, namely $\{\,A_0,\ A_1+\dots+A_{d-1},\ A_d\,\}$, and it is strongly regular.
\end{theorem}

\begin{corollary}\label{cor:C}
Every thick finite generalized octagon has $s\neq t$ (Feit--Higman). Hence the distance scheme of its point graph has no proper fusion. In particular the $5$-class scheme of Theorem~\ref{thm:A} has no proper fusion: apart from the trivial partition of $\{R_1,\dots,R_4\}$ into one block and the partition into singletons, no partition of the nontrivial orbitals into blocks gives an association scheme, and no strongly regular graph is a union of orbitals.
\end{corollary}

\paragraph{Relation to earlier work.}
The graph is not new. It is the collinearity graph of the Ree--Tits octagon~\cite{Tits83}; its description as the coset graph $\Gamma(T,H)$ with $\Aut=T{:}2$ appears in Crnkovi\'c--Rukavina--\v{S}vob~\cite[Thm.~9, Table~13]{CRS}. The array, the spectrum and the vanishing of a Krein parameter are routine consequences of the general theory of generalized polygons~\cite{BCN,CGS}. What we add is the explicit, independently verified orbital algebra of the group action (Theorem~\ref{thm:A}) and the classification of fusions (Theorem~\ref{thm:B}). We did not find a statement of Theorem~\ref{thm:B} in \cite{CRS}, \cite{BCN} or the papers citing \cite{CRS}; this is the outcome of a literature search and not a proof of novelty, and the result may well be known to specialists. The octagon case reduces, as Corollary~\ref{cor:C} shows, to the Feit--Higman theorem; the hexagon case with $s=t$ is the familiar fact that the split Cayley hexagon $H(q)$ sits inside the parabolic quadric $Q(6,q)$ (Remark~\ref{rem:hex}; see~\cite{VM}).

\section{Fusion criterion}

Let $(X,\{R_i\}_{i=0}^D)$ be a symmetric association scheme with adjacency matrices $A_i$, valencies $k_i$ and primitive idempotents $E_0=\tfrac1{|X|}J,E_1,\dots,E_D$. The eigenmatrix $P=(P_{ej})$ is defined by $A_j=\sum_e P_{ej}E_e$. A \emph{fusion} is a partition $\Lambda=\{B_1,\dots,B_r\}$ of $\{1,\dots,D\}$ such that $V_\Lambda=\operatorname{span}\{I,A_{B_1},\dots,A_{B_r}\}$, $A_B=\sum_{j\in B}A_j$, is closed under multiplication; then $(X,\{R_0,R_B\})$ is an $r$-class association scheme. It is \emph{proper} if $2\le r\le D-1$.

For $e,f\in\{1,\dots,D\}$ and $B\subseteq\{1,\dots,D\}$ put
\[
\delta^{ef}(B)=\sum_{j\in B}\bigl(P_{ej}-P_{fj}\bigr).
\]
Since $\sum_{j=0}^DA_j=J$ and $JE_e=0$ for $e\ge1$, we have $\sum_{j=1}^DP_{ej}=-1$ and therefore
\begin{equation}\label{eq:compl}
\delta^{ef}(B^c)=-\delta^{ef}(B).
\end{equation}

The following criterion is standard (cf.~\cite{BI}); we include a proof.\begin{lemma}[Fusion criterion]\label{lem:crit}
For a partition $\Lambda$ of $\{1,\dots,D\}$ into $r$ blocks let $e\sim_\Lambda f$ mean $\delta^{ef}(B)=0$ for all $B\in\Lambda$. This is an equivalence relation on $\{1,\dots,D\}$ with at least $r$ classes, and $\Lambda$ is a fusion if and only if it has exactly $r$ classes.
\end{lemma}
\begin{proof}
Write $A_B=\sum_e u_B(e)E_e$ with $u_B(e)=\sum_{j\in B}P_{ej}$. Let $\approx$ be the equivalence relation on $\{0,\dots,D\}$ with $e\approx f$ iff $u_B(e)=u_B(f)$ for all $B\in\Lambda$, and let $F\subseteq\C^{D+1}$ be the algebra of functions constant on $\approx$-classes; $F$ is identified with $\operatorname{span}\{\sum_{e\in K}E_e: K\ \text{a class}\}$. Then $V_\Lambda\subseteq F$ and $\dim V_\Lambda=r+1$, because the $A_B$ have disjoint supports. Since $J=I+\sum_BA_B\in V_\Lambda$ equals $|X|E_0$, the class of $0$ is $\{0\}$, so $\approx$ has $1+c$ classes, where $c$ is the number of classes of $\sim_\Lambda$ on $\{1,\dots,D\}$, and $1+c\ge r+1$.
If $c=r$ then $V_\Lambda=F$, which is closed under multiplication. Conversely suppose $V_\Lambda$ is closed under multiplication. If $e\not\approx f$ choose $x\in V_\Lambda$ with $x(e)\ne x(f)$. Polynomials in $x$ lie in $V_\Lambda$ (it contains $I$), so $V_\Lambda$ contains the indicator function of $\{g:x(g)=x(e)\}$; multiplying such indicators over finitely many separating elements gives the indicator of the $\approx$-class of $e$. Hence $F\subseteq V_\Lambda$, $V_\Lambda=F$ and $1+c=r+1$.
\end{proof}

\begin{lemma}\label{lem:count}
Let $\Lambda$ be a proper fusion.
\begin{enumerate}
\item[(a)] If $D=3$ then $r=2$, and $\Lambda=\{B,B^c\}$ is a fusion iff $\delta^{ef}(B)=0$ for at least one pair $\{e,f\}$.
\item[(b)] If $D=4$ and $r=2$, then $\Lambda=\{B,B^c\}$ is a fusion iff $\delta^{ef}(B)=0$ for at least two distinct pairs.
\item[(c)] If $D=4$ and $r=3$, then $\Lambda=\{\{a\},\{b\},\{c,d\}\}$ and a fusion requires a pair $\{e,f\}$ with $\delta^{ef}(\{a\})=\delta^{ef}(\{b\})=0$.
\end{enumerate}
\end{lemma}
\begin{proof}
By Lemma~\ref{lem:crit} and \eqref{eq:compl} the conditions $\delta^{ef}(B)=0$ for all blocks of a two-block partition reduce to one block. (a) The relation $\sim_\Lambda$ has at least $2$ classes on three points, so it has exactly $2$ iff some pair is related. (b) A two-class relation on four points has class sizes $(2,2)$ or $(3,1)$, hence at least two related pairs; conversely two related pairs give at most two classes. (c) Three classes on four points means one related pair $\{e,f\}$; then $\delta^{ef}(\{a\})=\delta^{ef}(\{b\})=0$, and $\delta^{ef}(\{c,d\})=0$ follows from \eqref{eq:compl}.
\end{proof}

\begin{lemma}\label{lem:A1}
If the scheme is metric (the point scheme of a distance-regular graph), the numbers $P_{e1}=\theta_e$ are pairwise distinct, so $\delta^{ef}(\{1\})\ne0$ for all $e\ne f$. In particular no proper fusion has $\{1\}$ (or $\{1\}^c$) as a block.
\end{lemma}
\begin{proof}
$A_1$ generates the Bose--Mesner algebra, which has dimension $D+1$, so $A_1$ has $D+1$ distinct eigenvalues.
\end{proof}

\section{Eigenmatrices of the point schemes}

Let $d\in\{3,4\}$ and let $\Gamma$ be distance-regular with $b_0=s(t+1)$, $b_i=st$ $(1\le i\le d-1)$, $c_i=1$ $(1\le i\le d-1)$, $c_d=t+1$, which is the array of a generalized $2d$-gon of order $(s,t)$ \cite[\S6.5]{BCN}. Thus $a_i=s-1$ for $1\le i\le d-1$, $a_d=(s-1)(t+1)$, $k_1=s(t+1)$, $k_i=s(t+1)(st)^{i-1}$ for $i<d$ and $k_d=s^dt^{d-1}$. Put
\[
m=\sqrt{st}\ (d=3),\qquad m=\sqrt{2st}\ (d=4).
\]

\begin{proposition}\label{prop:P}
The eigenvalues of $A_1$ are $s(t+1)$, $-t-1$, $s-1\pm m$ and, if $d=4$, also $s-1$. The eigenmatrix $P=(P_{ej})$, rows $e=0,\dots,d$ and columns $j=0,\dots,d$, is as follows.
\begin{align*}
d=3:\quad P&=\begin{pmatrix}
1 & s(t+1) & s^2t(t+1) & s^3t^2\\
1 & s-1+m & ms-m-s & -ms\\
1 & s-1-m & -ms+m-s & ms\\
1 & -t-1 & t(t+1) & -t^2
\end{pmatrix},\\[4pt]
d=4:\quad P&=\begin{pmatrix}
1 & s(t+1) & s^2t(t+1) & s^3t^2(t+1) & s^4t^3\\
1 & s-1+m & st+m(s-1)-s & st(s-1)-ms & -s^2t\\
1 & s-1 & -s(t+1) & -st(s-1) & s^2t\\
1 & s-1-m & st-m(s-1)-s & st(s-1)+ms & -s^2t\\
1 & -t-1 & t(t+1) & -t^2(t+1) & t^3
\end{pmatrix}.
\end{align*}
\end{proposition}
\begin{proof}
The quotient (intersection) matrix of $A_1$ is tridiagonal with entries $b_{j-1},a_j,c_{j+1}$; its characteristic polynomial is $(x-s(t+1))(x+t+1)(x^2-2(s-1)x+(s-1)^2-st)$ for $d=3$ and $(x-s(t+1))(x-(s-1))(x+t+1)(x^2-2(s-1)x+(s-1)^2-2st)$ for $d=4$. For each root $\theta$ the row $(P_{ej})_j$ is determined by $P_{e0}=1$, $P_{e1}=\theta$ and $\theta P_{ej}=b_{j-1}P_{e,j-1}+a_jP_{ej}+c_{j+1}P_{e,j+1}$. We verified symbolically, modulo $m^2=st$ resp.\ $m^2=2st$, that the displayed rows satisfy this recursion and the sum rule $\sum_jP_{ej}=0$ ($e\ge1$); the verification is reproduced by the accompanying script.
\end{proof}

\section{Hexagons}

\begin{proof}[Proof of Theorem~\ref{thm:B} for $d=3$]
Let $s,t\ge2$ and $m=\sqrt{st}>0$. By Lemma~\ref{lem:count}(a) we must decide, for $B\in\{\{1\},\{2\},\{3\}\}$, whether some pair $\{e,f\}\subset\{1,2,3\}$ has $\delta^{ef}(B)=0$. From Proposition~\ref{prop:P}:

\smallskip\noindent\emph{$B=\{1\}$.} Impossible by Lemma~\ref{lem:A1} (directly: the differences are $2m$, $s+t+m$, $s+t-m$, and $m<s+t$).

\smallskip\noindent\emph{$B=\{2\}$.} The column $j=2$ is $(ms-m-s,\ -ms+m-s,\ t(t+1))$. Then $\delta^{12}=2m(s-1)>0$ and $\delta^{23}=-m(s-1)-s-t(t+1)<0$. If $\delta^{13}=0$ then $m(s-1)=s+t^2+t$; squaring and using $m^2=st$ gives $f(s,t)=0$ with
\[
f(s,t)=t\,s^3-(2t+1)s^2-t(2t+1)s-t^2(t+1)^2 .
\]
This factors as $f(s,t)=(s^2+st+t^2)\bigl(st-(t+1)^2\bigr)$. The first factor is positive, so $s=(t+1)^2/t=t+2+1/t$, which is not an integer for $t\ge2$. (For real $t>0$ this point does satisfy $m(s-1)=s+t^2+t$, with $m=t+1$, so the condition is not vacuous over the reals.)

\smallskip\noindent\emph{$B=\{3\}$.} The column $j=3$ is $(-ms,\ ms,\ -t^2)$. Here $\delta^{12}=-2ms\ne0$ and $\delta^{23}=ms+t^2>0$, while $\delta^{13}=t^2-ms=0$ iff $t^4=m^2s^2=s^3t$ iff $s=t$.

So a proper fusion exists iff $s=t$, and then only for $\Lambda=\{\{1,2\},\{3\}\}$.

\smallskip\noindent\emph{Strong regularity.} For $s=t=q$ the eigenvalues of $A_1+A_2$ are, by Proposition~\ref{prop:P}, $k=q(q+1)(q^2+1)$ on $E_0$, $q^2-1$ on $E_1,E_3$ and $-q^2-1$ on $E_2$. A graph with exactly these eigenvalues is strongly regular with $\mu=k+rs'$ and $\lambda=\mu+r+s'$, where $r=q^2-1$, $s'=-q^2-1$, giving
\[
\Bigl(\tfrac{q^6-1}{q-1},\ q(q+1)(q^2+1),\ q^3+q^2+q-1,\ (q+1)(q^2+1)\Bigr).
\]
\end{proof}

\begin{remark}\label{rem:hex}
These are the parameters of the collinearity graph of $Q(6,q)$. For the split Cayley hexagon $H(q)$ the identification is classical: its points are all points of $Q(6,q)$ and two points are at distance $\le2$ in $H(q)$ iff they are collinear in $Q(6,q)$. The computation above shows that the parameter identity holds for any distance-regular graph with the array of a hexagon of order $(q,q)$. Direct computation of the intersection numbers for $q=2,3,4,5$ gave $(63,30,13,15)$, $(364,120,38,40)$, $(1365,340,83,85)$, $(3906,780,154,156)$.
\end{remark}

\section{Octagons}

Now $d=4$, $s,t\ge2$ are integers and $m=\sqrt{2st}>0$. By Proposition~\ref{prop:P} every difference $\delta^{ef}(B)$ is of the form $\alpha+\beta m$ with $\alpha,\beta\in\mathbb Z[s,t]$ (after reducing with $m^2=2st$); all $42$ cases (seven two-block partitions $\{B,B^c\}$, six pairs) are listed in Appendix~\ref{app:table}. The last column records the sign of $\alpha$ and $\beta$ for $s,t\ge2$ (obtained by substituting $s=2+u$, $t=2+v$ and inspecting the coefficients; $\ast$ means no sign is forced). We use the following consequences.

\begin{enumerate}
\item[(i)] \emph{Pair $\{1,3\}$.} Here $\alpha=0$ and $\beta\in\{2,\,2(s-1),\,-2s,\,0,\,2s,\,-2(s-1),\,2\}$ for $B=\{1\},\{2\},\{3\},\{4\},\{1,2\},\{1,3\},\{1,4\}$ respectively. For $s\ge2$ it vanishes only for $B=\{4\}$ (and its complement).
\item[(ii)] \emph{Pair $\{2,4\}$.} Here $\beta=0$ and $\delta=\alpha$. It vanishes only for: $B=\{3\}$ iff $s=t+1$; $B=\{4\}$ iff $s=t$; $B=\{1,3\}$ iff $t(s-t-1)=1$; $B=\{1,4\}$ iff $t(t-s)=1$. The last two force $t\mid1$, impossible for $t\ge2$.
\end{enumerate}

\begin{proof}[Proof of Theorem~\ref{thm:B} for $d=4$]
\emph{Step 1: two-block partitions} ($r=2$). By Lemma~\ref{lem:A1} the case $B=\{1\}$ is excluded. For each other $B$ we list the pairs that can vanish and show that at most one does, unless $B=\{4\}$ and $s=t$. We use that $s,t\ge2$ are integers and that $\alpha+\beta m=0$ with $\alpha,\beta$ of the same strict sign is impossible.

\smallskip\noindent$B=\{2\}$. Pairs $\{1,2\},\{1,3\},\{2,4\}$ cannot vanish, nor can $\{2,3\}$: $\alpha+\beta m=0$ means $m(s-1)=2st=m^2$, so $m=s-1$ and $(s-1)^2=2st$, whence $s\mid(s-1)^2$, $s=1$. Only $\{1,4\}$ and $\{3,4\}$ remain, and they have the same $\alpha$ and opposite $\beta=\pm(s-1)$; vanishing of both would force $\alpha=0$ and $(s-1)m=0$, contradicting $s\ge2$.

\smallskip\noindent$B=\{3\}$. $\{1,3\}$ is excluded by (i); $\{2,3\}$ and $\{3,4\}$ have $\alpha,\beta$ of equal strict sign; $\{1,2\}$ vanishes iff $m=2t(s-1)$, i.e.\ $s=2t(s-1)^2$, impossible for $s\ge2$, $t\ge2$. Only $\{1,4\}$ and $\{2,4\}$ remain. If both vanish then $s=t+1$ by (ii), so $m^2=2t(t+1)$ and $s^2-s+t^2+t=2t(t+1)=m^2$; the condition for $\{1,4\}$ becomes $m(tm-s)=0$, i.e.\ $m=(t+1)/t$, and then $(t+1)^2/t^2=2t(t+1)$ gives $t+1=2t^3$, false for $t\ge2$.

\smallskip\noindent$B=\{4\}$. All $\beta$ vanish; $\alpha$ is $-2s^2t$, $0$, $-t(s^2+t^2)$, $2s^2t$, $t(s-t)(s+t)$, $-t(s^2+t^2)$ for the pairs $\{1,2\},\{1,3\},\{1,4\},\{2,3\},\{2,4\},\{3,4\}$. So $\{1,3\}$ always vanishes, $\{2,4\}$ vanishes iff $s=t$, and no other pair does. By Lemma~\ref{lem:count}(b) the partition $\{\{1,2,3\},\{4\}\}$ is a fusion iff $s=t$.

\smallskip\noindent$B=\{1,2\}$. Pairs $\{1,2\},\{1,3\},\{2,4\}$ do not vanish. $\{2,3\}$ vanishes iff $m=2t$, i.e.\ $s=2t$. $\{1,4\}$ and $\{3,4\}$ have $\alpha=t(s-t)$, $\beta=\pm s$. $\{1,4\}$ needs $\alpha<0$, i.e.\ $t>s$, which is incompatible with $s=2t$; $\{3,4\}$ with $s=2t$ gives $\alpha+\beta m=t^2-4t^2\ne0$; and both cannot vanish (that would force $s=0$). So at most one pair vanishes.

\smallskip\noindent$B=\{1,3\}$. $\{1,2\}$ vanishes iff $m=2st$, hence $2st=4s^2t^2$, impossible; $\{1,3\}$, $\{2,3\}$, $\{3,4\}$ do not vanish by (i) and sign; $\{2,4\}$ is excluded by (ii). Only $\{1,4\}$ can vanish.

\smallskip\noindent$B=\{1,4\}$. $\{1,2\}$ vanishes iff $m=2s^2t$, hence $2s^3t=1$, impossible; $\{1,3\}$ by (i); $\{2,3\}$ and $\{3,4\}$ have $\alpha,\beta$ of equal strict sign; $\{2,4\}$ by (ii). Only $\{1,4\}$ can vanish.

Hence by Lemma~\ref{lem:count}(b) the only proper two-block fusion is $\{\{1,2,3\},\{4\}\}$, and it occurs iff $s=t$.

\smallskip\emph{Step 2: three-block partitions} ($r=3$). By Lemma~\ref{lem:count}(c) we need a pair $\{e,f\}$ and two distinct singletons $\{a\},\{b\}$ with $\delta^{ef}(\{a\})=\delta^{ef}(\{b\})=0$. By Lemma~\ref{lem:A1}, $a,b\in\{2,3,4\}$. Reading the singleton rows of Appendix~\ref{app:table} together with Step~1: in row $\{2\}$ only the pairs $\{1,4\}$ and $\{3,4\}$ can vanish; in row $\{3\}$ only $\{1,4\}$ and $\{2,4\}$; in row $\{4\}$ only $\{1,3\}$ and (if $s=t$) $\{2,4\}$. A pair must therefore occur in two rows. The pair $\{3,4\}$ occurs only in row $\{2\}$ and $\{1,3\}$ only in row $\{4\}$. The pair $\{2,4\}$ would need $s=t+1$ (row $\{3\}$) and $s=t$ (row $\{4\}$). The pair $\{1,4\}$ would need $\delta(\{2\})=\delta(\{3\})=0$, hence $\delta(\{2\})+\delta(\{3\})=0$, which reads
\[
m=t(s^2+t^2)-s-t .
\]
But for $s,t\ge2$ we have $t(s^2+t^2)\ge2(s^2+t^2)\ge(s+t)^2\ge4(s+t)$, so $m\ge3(s+t)$, whereas $m^2=2st\le(s+t)^2/2$ gives $m<s+t$. This contradiction shows that no three-block fusion exists.

\smallskip\emph{Step 3: strong regularity.} For $s=t$ the eigenvalues of $A_4$ are, by Proposition~\ref{prop:P}, $k_4=s^7$ on $E_0$, $-s^3$ on $E_1,E_3$ and $s^3$ on $E_2,E_4$, hence $A_4$ is strongly regular with parameters
$\bigl(\tfrac{s^8-1}{s-1},\ s^7,\ s^7-s^6,\ s^7-s^6\bigr)$.
\end{proof}

\begin{proof}[Proof of Corollary~\ref{cor:C}]
For a thick finite generalized octagon of order $(s,t)$ the Feit--Higman theorem~\cite{FH} requires $2st$ to be a perfect square, so $s=t$ is impossible ($2s^2$ is never a square). Theorem~\ref{thm:B} gives no proper fusion. A strongly regular graph that is a union $\bigcup_{i\in S}R_i$ of orbitals would give a two-block fusion $\{S,S^c\}$, which does not exist.
\end{proof}

\begin{remark}[Sharpness of thickness]\label{rem:sharp}
The hypothesis $s,t\ge2$ cannot be dropped. Direct computation from the intersection numbers (\texttt{thin/thin\_fusions.py}) shows that the thin octagon schemes of order $(1,2)$ and $(1,8)$ admit the fusions $\{\{2\},\{1,3\},\{4\}\}$ and $\{\{1,3\},\{2,4\}\}$, and that of order $(2,1)$ admits $\{\{2\},\{3\},\{1,4\}\}$ and $\{\{3\},\{1,2,4\}\}$ (these are the curve $s=1$ and the isolated point $(2,1)$ found by the exact computation described in Section~\ref{sec:scope}); the scheme of order $(8,1)$ has none. The fusions for order $(1,t)$ are imprimitive: the bipartition, and the halved graph, which is two copies of the collinearity graph of $\mathrm{GQ}(t,t)$. For order $(2,1)$ the graph of class $3$ is strongly regular with parameters $(45,16,8,4)$, those of the triangular graph $T(10)$. For $s=t$ the fusion $\{\{1,2,3\},\{4\}\}$ is formal: it would give strongly regular graphs such as $(255,128,64,64)$ for $s=2$ and $(3280,2187,1458,1458)$ for $s=3$, but no octagon of order $(s,s)$ exists.
\end{remark}

\section{The Tits group on $1755$ points}

\begin{proof}[Proof of Theorem~\ref{thm:A}]
Take $d=4$, $(s,t)=(2,4)$, so $m=4$. Then $|\Omega|=(1+s)(1+st)(1+s^2t^2)=3\cdot9\cdot65=1755$, the valencies are $k_i=1,10,80,640,1024$, and the array is $\{10,8,8,8;1,1,1,5\}$. The orbitals were computed from the ATLAS generators of the permutation representation of degree $1755$: the point stabilizer has orbits of lengths $1,10,80,640,1024$, a breadth-first search in the graph of the valency-$10$ orbital produces the layers $1,10,80,640,1024$ from a base point, and the $125$ intersection numbers $p_{ij}^k$ were obtained by counting, with a transporter element, $|R_i(x_0)\cap R_j(y)|$ for a representative $y\in R_k(x_0)$. These $125$ numbers agree in every entry with those derived from the intersection array by the recursion $A_1A_i=b_{i-1}A_{i-1}+a_iA_i+c_{i+1}A_{i+1}$. The distance relations of $\Gamma$ are $T$-invariant and there are five of them, equal to the rank, so each is a single orbital; in particular every orbital is symmetric and the algebra is commutative. The intersection matrices $B_i=(p_{ij}^k)_{j,k}$ ($B_0=I$) are
\begin{equation}\label{eq:B}
\begin{gathered}
B_1=\begin{pmatrix}0 & 1 & 0 & 0 & 0 \\ 10 & 1 & 1 & 0 & 0 \\ 0 & 8 & 1 & 1 & 0 \\ 0 & 0 & 8 & 1 & 5 \\ 0 & 0 & 0 & 8 & 5\end{pmatrix},\quad B_2=\begin{pmatrix}0 & 0 & 1 & 0 & 0 \\ 0 & 8 & 1 & 1 & 0 \\ 80 & 8 & 6 & 1 & 5 \\ 0 & 64 & 8 & 38 & 25 \\ 0 & 0 & 64 & 40 & 50\end{pmatrix} ,\\[6pt]
B_3=\begin{pmatrix}0 & 0 & 0 & 1 & 0 \\ 0 & 0 & 8 & 1 & 5 \\ 0 & 64 & 8 & 38 & 25 \\ 640 & 64 & 304 & 200 & 250 \\ 0 & 512 & 320 & 400 & 360\end{pmatrix},\quad B_4=\begin{pmatrix}0 & 0 & 0 & 0 & 1 \\ 0 & 0 & 0 & 8 & 5 \\ 0 & 0 & 64 & 40 & 50 \\ 0 & 512 & 320 & 400 & 360 \\ 1024 & 512 & 640 & 576 & 608\end{pmatrix}
\end{gathered}
\end{equation}
(rows $j$, columns $k$), and they satisfy $p_{ij}^k=p_{ji}^k$ and $\sum_kp_{ij}^kk_k=k_ik_j$. Proposition~\ref{prop:P} at $(s,t)=(2,4)$, $m=4$ gives the eigenmatrix
\begin{equation}\label{eq:P}
P=\begin{pmatrix}1 & 10 & 80 & 640 & 1024 \\ 1 & 5 & 10 & 0 & -16 \\ 1 & 1 & -10 & -8 & 16 \\ 1 & -3 & 2 & 16 & -16 \\ 1 & -5 & 20 & -80 & 64\end{pmatrix},\qquad \text{multiplicities }(m_0,\dots,m_4)=(1,351,650,675,78),
\end{equation}
with eigenvalues $10,5,1,-3,-5$ of $A_1$ and $m_e=|\Omega|\big/\sum_jP_{ej}^2/k_j$.
\end{proof}

\begin{remark}[Krein parameters]
With $Q$ determined by $QP=|\Omega|I$, all Krein parameters $q_{ef}^g$ are non-negative and the only vanishing one among the non-principal indices is $q_{44}^{4}=0$, where index $4$ is the eigenspace of dimension $78$ (eigenvalue $-5$). This is an instance of the extremal-polygon phenomenon of Cameron--Goethals--Seidel~\cite{CGS}. The scheme is not $Q$-polynomial with respect to any ordering of its eigenspaces. This is known: by De Bruyn--Vanhove~\cite{DBV}, for $d\ge4$ the only thick $Q$-polynomial regular near $2d$-gons are the dual polar spaces and the Hamming near polygons, and a thick generalized octagon is neither. We also checked it directly in exact arithmetic over all $24$ orderings (\texttt{qpoly/qpoly\_check.py}). The graph is not tight in the sense of Juri\v{s}i\'c--Koolen--Terwilliger~\cite{JKT}: with $k=10$, $a_1=1$, $b_1=8$, $\theta_1=5$, $\theta_4=-5$ the fundamental bound $(\theta_1+\tfrac{k}{a_1+1})(\theta_4+\tfrac{k}{a_1+1})\ge-\tfrac{ka_1b_1}{(a_1+1)^2}$ reads $0\ge-20$, which is strict; also $a_4=5\ne0$, while tight graphs have $a_d=0$ (\texttt{tits1755/tightness\_check.py}).
\end{remark}

\section{Scope and verification}\label{sec:scope}

Everything above was recomputed in exact arithmetic (rational numbers and symbolic algebra over $\mathbb Q(s,t,m)$) by the accompanying script \texttt{verify\_fusions.py}: the recursion check of Proposition~\ref{prop:P}; a direct test, from the intersection numbers alone, of the closure condition for every proper partition for all $841$ integer pairs $2\le s,t\le30$ and $d=3,4$, which found proper fusions exactly on the diagonal $s=t$ and only the fusion named in Theorem~\ref{thm:B}; the hexagon cubic; the numbers of Section~6.

An independent exact check, \texttt{fusion/certify\_fusions.py}, avoids eigenvalues and square roots altogether. It computes $p_{ij}^k$ as rational functions of $s,t$ from the intersection array, writes the closure conditions of every proper partition as polynomial equations in $\mathbb Q[s,t]$, and reduces them with a lexicographic Gr\"obner basis over $\mathbb Q$. Real solutions are isolated exactly. For all real $s,t\ge1$ the outcome is: for $d=3$, $\{\{1,2\},\{3\}\}$ closes iff $s=t$, and $\{\{1,3\},\{2\}\}$ closes iff $s=1$ or $st=(t+1)^2$; for $d=4$, $\{\{1,2,3\},\{4\}\}$ closes iff $s=t$, the partitions $\{\{1,3\},\{2,4\}\}$ and $\{\{1,3\},\{2\},\{4\}\}$ close iff $s=1$, the partitions $\{\{1,2,4\},\{3\}\}$ and $\{\{1,4\},\{2\},\{3\}\}$ close only at $(s,t)=(2,1)$, and no other proper partition closes. This agrees with Theorem~\ref{thm:B} and with Remark~\ref{rem:sharp}.

\emph{Provenance.} The group-side computations (stabilizer orbits, breadth-first search, transporter counts) were carried out in an earlier session of this project directly from the ATLAS generators and saved as a table of the $125$ numbers $p_{ij}^k$; in the present session we re-checked that this table agrees in all $125$ entries with the numbers derived from the intersection array, and that it satisfies $p_{ij}^k=p_{ji}^k$ and $\sum_kp_{ij}^kk_k=k_ik_j$.

The order $|T|=17971200$ was re-derived from the generators by Schreier--Sims. The five eigenspace characters were evaluated exactly on all $22$ classes and agree with irreducible characters of degrees $1,78,351,650,675$ in the GAP character table library. The permutation character $1_H^T$ is their sum, so it is multiplicity-free (\texttt{tits1755/gap/tits\_chars.g}). Uniqueness of $\mathrm{GO}(2,4)$ as a generalized octagon is not addressed. The fusion statements are about the distance scheme only. The orbital configurations of $T$ on cosets of its other maximal subgroups are listed in \texttt{tits1755/coherent\_configurations\_note.md}; their fusions are in \texttt{tits1755/coset\_fusions\_note.md}. The proofs of Theorem~\ref{thm:B} use $s,t\in\mathbb Z_{\ge2}$ (divisibility arguments) and the sign table of Appendix~\ref{app:table}, which was generated symbolically.

\appendix
\section{Pairwise differences for the octagon}\label{app:table}
For each two-block partition $\{B,B^c\}$ and each pair $\{e,f\}\subseteq\{1,2,3,4\}$, $\delta^{ef}(B)=\alpha+\beta m$ with $m^2=2st$. By \eqref{eq:compl} the entries for $B^c$ are the negatives. Last column: sign of $(\alpha,\beta)$ for $s,t\ge2$ ($\ast$: not sign-definite).

{\small
\begin{longtable}{@{}l c l l c@{}}
\toprule
$B$ & $(e,f)$ & $\alpha$ & $\beta$ & $\operatorname{sgn}(\alpha,\beta)$\\
\midrule
\endhead
$\{1\}$ & $(1,2)$ & $0$ & $1$ & $0,\,+$ \\
 & $(1,3)$ & $0$ & $2$ & $0,\,+$ \\
 & $(1,4)$ & $s + t$ & $1$ & $+,\,+$ \\
 & $(2,3)$ & $0$ & $1$ & $0,\,+$ \\
 & $(2,4)$ & $s + t$ & $0$ & $+,\,0$ \\
 & $(3,4)$ & $s + t$ & $-1$ & $+,\,-$ \\
\midrule
$\{2\}$ & $(1,2)$ & $2 s t$ & $s - 1$ & $+,\,+$ \\
 & $(1,3)$ & $0$ & $2 \left(s - 1\right)$ & $0,\,+$ \\
 & $(1,4)$ & $s t - s - t^{2} - t$ & $s - 1$ & $\ast,\,+$ \\
 & $(2,3)$ & $- 2 s t$ & $s - 1$ & $-,\,+$ \\
 & $(2,4)$ & $- \left(s + t\right) \left(t + 1\right)$ & $0$ & $-,\,0$ \\
 & $(3,4)$ & $s t - s - t^{2} - t$ & $1 - s$ & $\ast,\,-$ \\
\midrule
$\{3\}$ & $(1,2)$ & $2 s t \left(s - 1\right)$ & $- s$ & $+,\,-$ \\
 & $(1,3)$ & $0$ & $- 2 s$ & $0,\,-$ \\
 & $(1,4)$ & $t \left(s^{2} - s + t^{2} + t\right)$ & $- s$ & $+,\,-$ \\
 & $(2,3)$ & $- 2 s t \left(s - 1\right)$ & $- s$ & $-,\,-$ \\
 & $(2,4)$ & $- t \left(s + t\right) \left(s - t - 1\right)$ & $0$ & $\ast,\,0$ \\
 & $(3,4)$ & $t \left(s^{2} - s + t^{2} + t\right)$ & $s$ & $+,\,+$ \\
\midrule
$\{4\}$ & $(1,2)$ & $- 2 s^{2} t$ & $0$ & $-,\,0$ \\
 & $(1,3)$ & $0$ & $0$ & $0,\,0$ \\
 & $(1,4)$ & $- t \left(s^{2} + t^{2}\right)$ & $0$ & $-,\,0$ \\
 & $(2,3)$ & $2 s^{2} t$ & $0$ & $+,\,0$ \\
 & $(2,4)$ & $t \left(s - t\right) \left(s + t\right)$ & $0$ & $\ast,\,0$ \\
 & $(3,4)$ & $- t \left(s^{2} + t^{2}\right)$ & $0$ & $-,\,0$ \\
\midrule
$\{1,2\}$ & $(1,2)$ & $2 s t$ & $s$ & $+,\,+$ \\
 & $(1,3)$ & $0$ & $2 s$ & $0,\,+$ \\
 & $(1,4)$ & $t \left(s - t\right)$ & $s$ & $\ast,\,+$ \\
 & $(2,3)$ & $- 2 s t$ & $s$ & $-,\,+$ \\
 & $(2,4)$ & $- t \left(s + t\right)$ & $0$ & $-,\,0$ \\
 & $(3,4)$ & $t \left(s - t\right)$ & $- s$ & $\ast,\,-$ \\
\midrule
$\{1,3\}$ & $(1,2)$ & $2 s t \left(s - 1\right)$ & $1 - s$ & $+,\,-$ \\
 & $(1,3)$ & $0$ & $- 2 \left(s - 1\right)$ & $0,\,-$ \\
 & $(1,4)$ & $s^{2} t - s t + s + t^{3} + t^{2} + t$ & $1 - s$ & $+,\,-$ \\
 & $(2,3)$ & $- 2 s t \left(s - 1\right)$ & $1 - s$ & $-,\,-$ \\
 & $(2,4)$ & $- \left(s + t\right) \left(s t - t^{2} - t - 1\right)$ & $0$ & $\ast,\,0$ \\
 & $(3,4)$ & $s^{2} t - s t + s + t^{3} + t^{2} + t$ & $s - 1$ & $+,\,+$ \\
\midrule
$\{1,4\}$ & $(1,2)$ & $- 2 s^{2} t$ & $1$ & $-,\,+$ \\
 & $(1,3)$ & $0$ & $2$ & $0,\,+$ \\
 & $(1,4)$ & $- s^{2} t + s - t^{3} + t$ & $1$ & $-,\,+$ \\
 & $(2,3)$ & $2 s^{2} t$ & $1$ & $+,\,+$ \\
 & $(2,4)$ & $\left(s + t\right) \left(s t - t^{2} + 1\right)$ & $0$ & $\ast,\,0$ \\
 & $(3,4)$ & $- s^{2} t + s - t^{3} + t$ & $-1$ & $-,\,-$ \\
\bottomrule
\end{longtable}
}

\begin{thebibliography}{99}
\bibitem{ATLASrep} R.~A.~Wilson et al., \emph{ATLAS of Group Representations}, version 3 (online): permutation representation \texttt{TF42G1-p1755B0} of ${}^2F_4(2)'$.
\bibitem{BCN} A.~E.~Brouwer, A.~M.~Cohen, A.~Neumaier, \emph{Distance-Regular Graphs}, Springer, Berlin, 1989, doi:10.1007/978-3-642-74341-2.
\bibitem{BI} E.~Bannai, T.~Ito, \emph{Algebraic Combinatorics I: Association Schemes}, Mathematics Lecture Note Series 58, Benjamin/Cummings, Menlo Park, CA, 1984, ISBN 0-8053-0490-8.
\bibitem{CGS} P.~J.~Cameron, J.-M.~Goethals, J.~J.~Seidel, Strongly regular graphs having strongly regular subconstituents, \emph{J.\ Algebra} \textbf{55} (1978) 257--280.
\bibitem{CRS} D.~Crnkovi\'c, S.~Rukavina, A.~\v{S}vob, On some distance-regular graphs with many vertices, \emph{J.\ Algebraic Combin.} \textbf{51} (2020) 641--652, doi:10.1007/s10801-019-00888-5.
\bibitem{DBV} B.~De Bruyn, F.~Vanhove, On $Q$-polynomial regular near $2d$-gons, \emph{Combinatorica} \textbf{35} (2015) 181--208, doi:10.1007/s00493-014-3039-x.
\bibitem{FH} W.~Feit, G.~Higman, The nonexistence of certain generalized polygons, \emph{J.\ Algebra} \textbf{1} (1964) 114--131.
\bibitem{VM} H.~Van Maldeghem, \emph{Generalized Polygons}, Birkh\"auser, 1998.
\bibitem{JKT} A.~Juri\v{s}i\'c, J.~Koolen, P.~Terwilliger, Tight distance-regular graphs, \emph{J.\ Algebraic Combin.} \textbf{12} (2000) 163--197, doi:10.1023/A:1026544111089.
\bibitem{Tits83} J.~Tits, Moufang octagons and the Ree groups of type ${}^2F_4$, \emph{Amer.\ J.\ Math.} \textbf{105} (1983) 539--594, doi:10.2307/2374268.
\end{thebibliography}

\end{document}
